\section{Tổng kết}

\begin{frame}{\secname: Tổng kết}
\begin{itemize}
    \item Báo cáo đã tìm hiểu, phân tích và cài đặt LViT cho bài toán phân đoạn ảnh y khoa theo câu nhắc văn bản.
    \item LViT kết hợp đặc trưng ảnh và ngôn ngữ, giúp mô hình định hướng vùng tổn thương theo mô tả lâm sàng.
    \item Cách tiếp cận này phù hợp với dữ liệu y khoa, nơi ảnh thường đi kèm ghi chú vị trí, mức độ hoặc đặc điểm bệnh lý.
\end{itemize}
\end{frame}

\begin{frame}{\secname: Tổng kết}
\begin{itemize}
    \item Báo cáo đã phân tích các thành phần chính: CNN U-shape, ViT U-shape, BERT embedding, Patch Embedding, Reconstruct và PLAM.
    \item Nhánh CNN giữ chi tiết cục bộ, trong khi nhánh ViT học quan hệ toàn cục giữa các vùng ảnh.
    \item Văn bản được dùng như tín hiệu điều kiện, giúp quá trình phân đoạn không chỉ phụ thuộc vào ảnh đầu vào.
\end{itemize}
\end{frame}

\begin{frame}{\secname: Tổng kết}
\begin{itemize}
    \item Báo cáo đã trình bày pipeline huấn luyện, suy luận và tính loss kết hợp với sinh nhãn giả.
    \item Language-Vision Loss và EPI cho thấy tiềm năng khai thác dữ liệu ít nhãn hoặc không nhãn.
    \item Phần cài đặt đã được điều chỉnh để tương thích với môi trường Kaggle, dữ liệu thực nghiệm và tài nguyên giới hạn.
\end{itemize}
\end{frame}

\section{Hạn chế}

\begin{frame}{\secname: Hạn chế}
\begin{itemize}
    \item Thực nghiệm chủ yếu tập trung trên QaTa-COV19-v2 và MosMedData+, chưa đánh giá đầy đủ trên nhiều miền dữ liệu khác.
    \item Cơ chế bán giám sát và sinh nhãn giả EPI mới được tái hiện ở mức hạn chế, chưa tương đương hoàn toàn với công trình gốc.
    \item Tập ESO-CT không được sử dụng do chưa được công bố công khai, làm giảm phạm vi đánh giá.
\end{itemize}
\end{frame}

\begin{frame}{\secname: Hạn chế}
\begin{itemize}
    \item Kết quả tái lập chịu ảnh hưởng bởi khác biệt môi trường phần mềm, cấu trúc dữ liệu và cách chia tập.
    \item Giới hạn tài nguyên trên Kaggle khiến quá trình huấn luyện khó tái hiện đầy đủ điều kiện của công trình gốc.
    \item Một số siêu tham số và bước tiền xử lý phải điều chỉnh, nên kết quả có thể khác so với kết quả công bố.
\end{itemize}
\end{frame}

\section{Hướng phát triển}

\begin{frame}{\secname: Hướng phát triển}
\begin{itemize}
    \item Hoàn thiện thực nghiệm bán giám sát bằng cách triển khai đầy đủ cơ chế sinh và cập nhật nhãn giả EPI.
    \item Mở rộng đánh giá trên nhiều loại ảnh y khoa hơn như CT, MRI, X-ray và ảnh mô bệnh học.
    \item Kiểm tra khả năng tổng quát hóa của LViT trên nhiều bệnh lý, nhiều cách chia dữ liệu và nhiều dạng chú thích văn bản.
\end{itemize}
\end{frame}

\begin{frame}{\secname: Hướng phát triển}
\begin{itemize}
    \item Nghiên cứu các bộ mã hóa văn bản mạnh hơn để khai thác tốt hơn ngữ cảnh trong mô tả y khoa.
    \item Tích hợp các mô hình nền tảng y khoa nhằm cải thiện căn chỉnh giữa đặc trưng ảnh và đặc trưng ngôn ngữ.
    \item Tối ưu quy trình sinh, lọc và cập nhật nhãn giả để giảm chi phí chú thích dữ liệu y khoa.
\end{itemize}
\end{frame}

\begin{frame}{\secname: Kết luận chung}
\begin{itemize}
    \item LViT là hướng tiếp cận tiềm năng cho phân đoạn ảnh y khoa đa phương thức.
    \item Việc kết hợp thị giác và ngôn ngữ giúp tận dụng tri thức mô tả từ chuyên gia.
    \item Dù còn hạn chế trong tái lập, báo cáo đã làm rõ kiến trúc, cơ sở phương pháp và các vấn đề kỹ thuật khi triển khai LViT.
\end{itemize}
\end{frame}